% from ega4-16-fr.tex
\oldpage[IV-4]{49}
inverse à droite (16.1.7). On peut donc supposer
$X=\operatorname{Spec}(B)$ et $Y=\operatorname{Spec}(A)$ affines, $B$ étant
une $A$-algèbre augmentée, et l'idéal d'augmentation $\mathfrak{I}$ étant de
type fini. Comme $\mathscr{A}^{(n)}$ s'identifie alors à
$(B/\mathfrak{I}^{n+1})^{\sim}$, le corollaire résulte de l'équivalence de b)
et c) dans (0, 19.5.4) puisque $B/\mathfrak{I}=A$.

On notera que, dans le cas affine envisagé, le fait que $j$ soit une immersion
quasi-régulière équivaut encore, en vertu de (0, 19.5.4), à dire que $B$ est
une $A$-algèbre formellement lisse pour la topologie
$\mathfrak{I}$-pré-adique.

On notera aussi que la condition que $j$ soit une immersion localement de
présentation finie est toujours vérifiée lorsque le morphisme $X\to Y$ est
localement de type fini (1.4.3, (v)).

\begin{proposition}[16.9.10]
\EGAReaderTarget{ega4.c04.prop.16-9-10}{IV.16.9.10}
Soient $X$ un préschéma localement noethérien, $Y$ un sous-préschéma de $X$,
$j:Y\to X$ l'injection canonique, $y$ un point de $Y$.
\begin{enumerate}
  \item[(i)] Pour qu'il existe un voisinage ouvert $U$ de $y$ dans $X$ tel que
  la restriction $Y\cap U\to U$ de $j$ soit une immersion régulière, il faut
  et il suffit que le noyau $\mathscr{I}_y$ de l'homomorphisme surjectif
  $\mathcal{O}_{X,y}\to\mathcal{O}_{Y,y}$ soit engendré par une suite régulière
  d'éléments de $\mathcal{O}_{X,y}$.
  \item[(ii)] Pour que l'immersion $j$ soit régulière, il faut et il suffit
  qu'elle soit quasi-régulière.
\end{enumerate}
\end{proposition}

\begin{enumerate}
  \item[(i)] On peut se borner au cas où $Y$ est un sous-préschéma fermé de
  $X$ défini par un Idéal \emph{cohérent} $\mathscr{I}$ de $\mathcal{O}_X$.
  La condition est évidemment nécessaire. Inversement, si $\mathscr{I}_y$
  est engendré par une suite régulière $(s_i)_y$, où les $s_i$ sont des
  sections de $\mathscr{I}$ au-dessus d'un voisinage ouvert $U$ de $y$ dans
  $X$, on peut supposer que les $s_i$ engendrent $\mathscr{I}|U$
  (0\textsubscript{I}, 5.2.2) et forment une suite régulière (0, 15.2.4),
  d'où l'assertion.
  \item[(ii)] Le fait qu'une immersion quasi-régulière soit régulière résulte
  de (i) et de l'identité des suites quasi-régulières et des suites régulières
  dans $\mathcal{O}_{X,y}$, formées d'éléments de l'idéal maximal
  (0, 15.1.11).
\end{enumerate}

Lorsque (sans hypothèse noethérienne sur $X$) le noyau $\mathscr{I}_y$ de
$\mathcal{O}_{X,y}\to\mathcal{O}_{Y,y}$ est engendré par une suite régulière
d'éléments de $\mathcal{O}_{X,y}$, on dit que l'immersion $j$ est
\emph{régulière au point $y$}.

\begin{corollary}[16.9.11]
\label{IV.16.9.11}
Soit $X$ un préschéma localement noethérien; alors tout Idéal quasi-régulier de
$\mathcal{O}_X$ est régulier.
\end{corollary}

\begin{remarks}[16.9.12]
\label{IV.16.9.12}
\begin{enumerate}
  \item[(i)] On notera qu'une immersion régulière n'est pas en général un
  morphisme plat, ni \emph{a fortiori} un morphisme \emph{régulier} au sens de
  (6.8.1).
  \item[(ii)] Soit $A$ un anneau local noethérien; il résulte aussitôt de
  (16.9.4) et de (0, 17.1.1) que pour que $A$ soit \emph{régulier}, il faut et
  il suffit que son idéal maximal $\mathfrak{m}$ soit \emph{quasi-régulier}
  (ou \emph{régulier}, ce qui revient au même puisque $A$ est noethérien).
  Pour qu'un schéma affine noethérien $X$ soit \emph{régulier}, il faut et il
  suffit que pour tout point fermé $x\in X$, l'injection canonique
  $\operatorname{Spec}(k(x))\to X$ soit une immersion \emph{régulière}.
\end{enumerate}
\end{remarks}

\begin{proposition}[16.9.13]
\label{IV.16.9.13}
Soient $X$ un préschéma localement noethérien, $Y$ un sous-préschéma de $X$,
$Y'$ un sous-préschéma de $Y$, tel que l'injection canonique $j:Y'\to Y$ soit
régulière. Alors la suite de $\mathcal{O}_{Y'}$-Modules
\[
  \xymatrix{
    0\ar[r] & j^*(\mathcal{N}_{Y/X})\ar[r]
      & \mathcal{N}_{Y'/X}\ar[r] & \mathcal{N}_{Y'/Y}\ar[r] & 0
  }
  \tag{16.9.13.1}
  \label{IV.16.9.13.1}
\]
